\scanpage{27}
% archive sheet number: 26
% L1--4
donc \deleted{\unsure{il s’ensuit}} \inserted{\unsure{il s’ensuit}} \deleted{\unsure{fac.}} \unsure{facile} que les $\unsure{\mathcal I|\mathscr I|\Xi}^{\unsure{n|m}}$ sont tels que $\unsure{\underline{\mathcal O}|\underline O}_X/\mathcal I^n$ sont encore plats sur $Y$. Soit $\deleted{\unsure{Y}}S_n$ le sous-schéma de $X$ défini par $\mathcal I^n$, et soit $X_n=X\amalg_{S_n}Y$ le $\deleted{S}Y$-schéma obtenu en ``\unsure{contractant} $S_n$ en $S$''.
% L5--6
\deleted{\unsure{Soit}} \deleted{\unsure{de sorte} $X_{\unsure{S|\xi}}=X-S$.} Donc $X_0=X$, et les $X_i$ forment un \uline{système} \uline{\deleted{projectif}\inserted{inductif}}.
% left, cancelled square and replacement square
\diagram{\begin{tikzcd}[ampersand replacement=\&]
\deleted{\unsure{S}_n}\arrow[r,no head]\arrow[d]\arrow[dr,no head,"\deleted{\unsure{/}}" description] \& \deleted X\arrow[d] \& S_n\arrow[r,"i_n"]\arrow[d,"f i_n"'] \& X\arrow[d,"p_n"]\\
\deleted Y\arrow[r,no head] \& \deleted{X_n} \& Y\arrow[r,"\varepsilon_n"'] \& X_n
\end{tikzcd}}
\diagram{\begin{tikzcd}[ampersand replacement=\&]
\cdots\& X_2\arrow[l,no head]\&X_1\arrow[l]\&X_0=X\arrow[l]
\end{tikzcd}}
% L7--14
Donc les $\unsure{\mathscr P|\mathcal P}_{X_n/S}$ forment un système \deleted{\uline{inductif}} projectif, d’où un pro-($S$-préschéma en groupes) noté $\mathscr P_{S,X/Y}$, et qui ne dépend que de la \uline{partie fermée} $S$ de $X/Y$, comme on constate facilement (du moins, si $Y$ est noethérien) --- cel\unsure{a} résulte d’ailleurs \unsure{de} la ``construction'' qu’on en donne plus bas.

% L15--17
Les morphismes canoniques $X-S\longrightarrow\mathscr P^1_{X_n/Y}$ définissent à la limite un morphisme
\[
\boxed{\unsure{\varphi|\phi}_S:X-S\longrightarrow\mathscr P^1_{S,X/Y}\subset\mathscr P_{S,X/Y}}.
\]
